Sources : Eurostat, gov\_10a\_taxag catégorie D211 et nama\_10\_gdp poste P31\_S14\_S15, données 2024 ; taux normaux de janvier 2026 relevés par la Tax Foundation.\par Note : l'écart à 100 \% mesure ensemble les taux réduits, les exonérations et ce qui échappe. Il ne les sépare pas. La comparaison reste significative parce que le dénominateur est le même pour tous les pays, mais le niveau dépend de la définition retenue de la consommation : élargie à la consommation publique, le ratio français tomberait autour de la moitié.
